\scanpage{103}
But we still have to make explicit the structure of $S\Pi M_G$, in terms of the $M_{g_\alpha,\widehat I_\alpha}$ — we shall do so in the next number.

For an arbitrary normalization morphism \unsure{(}not necessarily strict), the commutative diagram gives a commutative diagram (up to isomorphism of functors)
\diagram{\begin{tikzcd}[
ampersand replacement=\&, column sep=small, row sep=large,
execute at end picture={
  \node[left=2em] at (\tikzcdmatrixname-1-1.west) {$(146)$};
  \path (\tikzcdmatrixname-1-1.east) ++(.35,0) coordinate (topstart);
  \path (\tikzcdmatrixname-1-3.west) ++(-.35,0) coordinate (topfinish);
  \draw[->] (topstart) -- (topfinish |- topstart);
  \draw (topstart) -- ++(.20,-.04);
  % The short lower-left return is visible; its interpretation as a head is uncertain.
  \node[below=1pt,inner sep=0pt] at (topstart) {$\unsure{\vphantom{-}}$};
}]
\Pi_1M_{\substack{\deleted{\unsure{G'}}\\G}} \arrow[d,"\wr"]
\&\& \Pi_1M_{G'} \arrow[d,"\wr"] \\
\Pi_1M_{\widetilde G}=\prod_{\alpha\in S}M_{g_\alpha,\widehat I_\alpha}
\arrow[rr]
\arrow[dr,"{\substack{\text{we identify }S'\\\text{with a subset of}\\S,\ \text{and }g_\alpha\text{ with }g'_{\alpha'}}}"']
\&\&
\Pi_1M_{\widetilde G'}=\prod_{\alpha'\in S'}\raisebox{-0.7ex}{\scriptsize$\deleted{\unsure{1}}$}M_{g'_{\alpha'},\widehat I'_{\unsure{\alpha'|\alpha}}} \\
\& \prod_{\alpha'\in S'}\raisebox{-0.7ex}{\scriptsize$\deleted{\unsure{*|1}}$}M_{g_\alpha,\widehat I_\alpha}
\arrow[ur,"\psi"'] \&
\end{tikzcd}}
\margin{This comes from an analogous diagram at the level of the modular multiplicities themselves\ldots}

The oblique arrow $\psi$ is induced by arrows
\[
(147)\qquad \Pi_1M_{g_\alpha,\widehat I_\alpha}\longrightarrow\Pi_1M_{g_\alpha^{\deleted{\unsure{'}}},\widehat I'_\alpha}
\]
in view of the fact that we have $\widehat I'_\alpha\subset\widehat I_\alpha$.

Thus knowledge of \inserted{general} normalization morphisms is reduced to that of the morphisms
\diagram{\begin{tikzcd}[
ampersand replacement=\&, column sep=2.2em,
execute at end picture={
  \node[left=2.2em] at (\tikzcdmatrixname-1-1.west) {$(148)$};
  \path (\tikzcdmatrixname-1-1.east) ++(.12,0) coordinate (tail148);
  \path (\tikzcdmatrixname-1-2.west) ++(.31,0) coordinate (overlap148);
  \path (overlap148 |- tail148) coordinate (end148);
  \path (tail148) -- (end148) coordinate[pos=.46] (head148);
  % One continuous shaft bears two right heads; its second head crosses the target Pi.
  \draw (tail148) -- (end148);
  \draw (head148) ++(-.10,.05) -- (head148) -- ++(-.10,-.05);
  \draw (end148) ++(-.10,.05) -- (end148) -- ++(-.10,-.05);
}]
\Pi_1M_{g,I} \& \Pi_1M_{g,I'} \& \text{for }I'\subset I,
\end{tikzcd}}
\margin{102}
